\begin{table}
\begin{tcolorbox}[tabularx={X|p{2cm}|X|},enhanced,fonttitle=\bfseries\large,fontupper=\normalsize\sffamily,
colback=yellow!10!white,colframe=red!50!black,colbacktitle=blue!30!white,
coltitle=black,title=Creating Graphs (Part 2)]
\hline
Command & Arguments & Purpose \\
\hline
\hline
\verb'-subset' & label  labeltex subset  & 
Define vertex coloring with two colors based on a subset of vertices. \\
\hline
\verb'-disjoint_sets_graph' & fname  & 
Reads a set of sets from the given file.
Define a graph on the sets. 
Vertices correspond to the given sets. 
Two vertices are adjacent 
if the associated sets are disjoint. \\
\hline
\verb'-orbital_graph' & $G$ $i$    & 
Define orbital graph from 
the $i$-th orbit of the group $G$ acting on pairs.  \\
\hline
\verb'-collinearity_graph' & inc-matrix   & 
Collinearity graph of the given incidence matrix.\\
\hline
\verb'-chain_graph' & P1 P2   & 
Chain graph with respect to the partitions P1 and P2.\\
\hline
\verb'-Neumaier_graph_16' &    & 
Create a Neumaier graph on 16 vertices.\\
\hline
\verb'-Neumaier_graph_25' &    & 
Create a Neumaier graph on 25 vertices, 
see~\cite{AbiadDeBoeckZeijlemaker2024}\\
\hline
\verb'-adjacency_bitvector' & bitvec N   & 
Create a graph on $N$ vertices from a bitvector.\\
\hline
\verb'-Cayley_graph' & $G$ gens   & 
Cayley graph with respect to group $G$ and generating set gens.\\
\hline
\end{tcolorbox}
\caption{\label{tab:graphtypes:two}Creating Graphs (Part 2)}
\end{table}
